\section{Step 81 Admit the level and writer windows}\label{h263-coeff-s01}
Validate request, level and context alignment and non-wrapping extents. Check the
block-type selector, reserved zero, writer capacity and existing bit count before
reading levels. The writer's pointer and capacity will remain unchanged; retain
the initial count for later complete-plan admission and commit.

\paragraph{Observe} Natural-order level input needs only halfword alignment.
A wide malformed selector is rejected before any level or buffer access.
\Needspace{8\baselineskip}
\paragraph{Check} What representation enters this function?
\begin{itemize}
  \item[A.] Selected coefficient levels, with separate INTRA DC code.
  \item[B.] Unquantized forward coefficients.
  \item[C.] Raw RGB pixels.
\end{itemize}

\section{Step 82 Validate every level before publication}\label{h263-coeff-s02}
Load signed halfwords for ordinary levels and check $-127\leq L\leq127$.
For INTRA index zero, validate the positive DC code and its excluded value 128.
Complete this pass even if an early event could already be encoded. A forbidden
last value must not leave earlier events in the caller stream.

\paragraph{Observe} Ordinary minus 128 is invalid. INTRA code 255 is valid and
means reconstructed DC 1024 at the later decoding boundary.
\Needspace{8\baselineskip}
\paragraph{Check} Which ordinary escape level is forbidden?
\begin{itemize}
  \item[A.] Minus 127.
  \item[B.] Minus 128.
  \item[C.] Plus one.
\end{itemize}

\section{Step 83 Locate the final nonzero in scan order}\label{h263-coeff-s03}
Use the scan permutation to inspect levels in reverse scan order until finding
a nonzero ordinary value. Stop at position zero for INTER or one for INTRA.
If none exists, retain an empty ordinary-event sequence. The INTRA DC code is
handled separately and must not become an AC event.

\paragraph{Observe} The default final nonzero is scan position two, natural
index eight. The remaining AC zeros are implicit.
\Needspace{8\baselineskip}
\paragraph{Check} What does LAST one describe?
\begin{itemize}
  \item[A.] The final allocated byte in the buffer.
  \item[B.] Always scan position sixty-three.
  \item[C.] The final nonzero ordinary event.
\end{itemize}

\section{Step 84 Count zeros between nonzero events}\label{h263-coeff-s04}
Walk forward from the first ordinary scan position. Increment RUN on zero levels.
For a nonzero, form an event with the current run and set LAST according to the
final position found earlier. Then reset RUN to zero. Address each level through
the supplied natural index, multiplied by two for halfword storage.

\paragraph{Observe} Two consecutive nonzeros both have RUN zero, even when their
natural indices are one and eight rather than adjacent memory locations.
\Needspace{8\baselineskip}
\paragraph{Check} When does the run counter reset?
\begin{itemize}
  \item[A.] After emitting each nonzero event.
  \item[B.] Only at the end of a picture.
  \item[C.] After every zero level.
\end{itemize}

\section{Step 85 Build a valid escape record first}\label{h263-coeff-s05}
Pack escape prefix three in seven bits, LAST in one, RUN in six and the validated
level's low eight bits. These fields total twenty-two bits and fit a single call
to the existing bit writer. Retain the full declared width so the prefix's leading
zeros are emitted.

For the first implementation, escape events can represent every legal ordinary
combination. This gives a complete baseline path before adding compact lookup.
Keep the separate INTRA DC field as eight bits without the escape prefix.

\paragraph{Observe} A negative escape level uses signed eight-bit representation;
it does not use a compact code's separate trailing sign convention.
\Needspace{8\baselineskip}
\paragraph{Check} How is minus 127 stored in an escape LEVEL field?
\begin{itemize}
  \item[A.] As a sign bit followed by a seven-bit VLC prefix.
  \item[B.] As the eight bits \texttt{10000001}.
  \item[C.] As forbidden level zero.
\end{itemize}

\section{Step 86 Select compact codes when available}\label{h263-coeff-s06}
Look up LAST, RUN and the positive magnitude in the supplied complete table.
When found, use its code and declared width, setting the final sign bit for a
negative level. Otherwise retain the escape record. Keep both implementations
for correctness and later size/work comparison.

A code's integer value does not encode its width. The compact final $+1$ code
has leading zero, so passing only its nonzero bit length would change the syntax.

\paragraph{Observe} The default compact block uses sixteen bits; the escape-only
block uses fifty-two. Both describe the same DC and AC levels.
\Needspace{8\baselineskip}
\paragraph{Check} What identifies one compact table entry?
\begin{itemize}
  \item[A.] Only the sign of the level.
  \item[B.] The current raw pixel value.
  \item[C.] LAST, RUN and level magnitude.
\end{itemize}

\section{Step 87 Admit the complete chosen plan}\label{h263-coeff-s07}
Store each code value and width privately and sum the widths. Include the INTRA
DC field if present. Compare the total with the caller's remaining bit capacity
before loading output bytes. Return status two when this selected representation
does not fit. An empty INTER plan succeeds without output access.

A compact plan can fit a small buffer where an escape-only plan cannot. Both
remain valid implementations. The public success cases supply enough space for
the worst-case legal representation, except for explicit fixed-size boundary tests.

\paragraph{Observe} A lone final INTER level minus 127 requires a twenty-two-bit
escape. At starting bit two, exactly three destination bytes suffice.
\Needspace{8\baselineskip}
\paragraph{Check} What must be included in capacity preflight?
\begin{itemize}
  \item[A.] Every field in the complete selected block plan.
  \item[B.] Only its first event.
  \item[C.] Only fields whose numerical values are nonzero.
\end{itemize}

\section{Step 88 Emit into a private byte window}\label{h263-coeff-s08}
Copy the touched destination bytes into private storage. Initialize a private
writer context at the caller's bit offset modulo eight. Emit every planned field
through your completed bit writer and check each child status. If a child fails,
return status three without any caller output or counter store.

After all children succeed, copy back only touched bytes, retaining their old
prefix/suffix bits, then commit the caller's total bit count. Save and restore
the return address across nested calls.

\paragraph{Observe} A late child can fail after earlier private fields were written.
The old public stream must still be intact.
\Needspace{8\baselineskip}
\paragraph{Check} Why is restoring only the old public count insufficient?
\begin{itemize}
  \item[A.] The count is a pixel pointer.
  \item[B.] Earlier public appends may already have replaced buffer bits.
  \item[C.] LAST fields are always native doublewords.
\end{itemize}

\section{Step 89 Decode emitted bits independently}\label{h263-coeff-s09}
Run public feedback on your completed source:
\begin{verbatim}
python3 drafts/h263-coefficient-construction/grade.py \
    my-coefficients.asm --cli /path/to/mars-cli.jar \
    --report coefficient-results.json
\end{verbatim}
The independent decoder recognizes all compact codes and the baseline escape.
It reconstructs the natural-order level array, rejects forbidden levels, excessive
runs, missing LAST and trailing bits, and compares every recovered input level.
The checker also verifies committed length, prefix/suffix preservation, guarded
state and access boundaries in both native byte orders.

\paragraph{Observe} Equivalent legal encodings can have different bit counts
without losing any selected coefficient information.
\Needspace{8\baselineskip}
\paragraph{Check} Which comparison establishes this block payload's meaning?
\begin{itemize}
  \item[A.] Matching the reference's private buffer bytes exactly.
  \item[B.] Matching only the first emitted byte.
  \item[C.] Decoding the payload and recovering every selected input level.
\end{itemize}

\section{Step 90 Keep block payload and macroblock flags separate}\label{h263-coeff-s10}
Retain the distinction between an empty ordinary payload and an INTRA DC field.
The next macroblock stage will select presence flags for its four Y and two chroma
blocks, write motion/type information, and invoke block emission in the required
order. This entry emits none of those flags and adds no arbitrary alignment.

Compare compact and escape payload sizes on the same selected levels. If measuring
execution work, include scanning, lookup, staging, child calls and publication.
Neither fewer bits nor one correct block proves whole-stream interoperability,
reference-picture ownership or complete-encoder speed.

\paragraph{Observe} INTRA with no nonzero AC still emits its DC byte. INTER with
no nonzero levels emits no block payload here.
\Needspace{8\baselineskip}
\paragraph{Check} Which later layer decides whether a block has ordinary payload?
\begin{itemize}
  \item[A.] The macroblock's coded-block presence information.
  \item[B.] Unused bits left in the output allocation.
  \item[C.] The host machine's native endianness.
\end{itemize}
